\section{Introduction}
\label{sec:intro}

Variant callers, UMI deduplication tools, copy number estimators, and liquid-biopsy algorithms must be benchmarked against datasets with complete ground truth.
Real patient data cannot provide this.
Synthetic simulation fills the gap.

Cancer sequencing introduces challenges that generic read simulators do not address.
Tumour samples are mixtures of cancer and normal cells, parametrised by purity and clonal architecture.
Allele frequencies are heterogeneous due to sub-clonal evolution.
Copy number alterations affect both local coverage and effective VAF.
Library preparation introduces artefacts specific to FFPE tissue and oxidative damage.
Cell-free DNA from liquid biopsy has a distinct fragment length distribution that read simulators do not model.

No single existing tool covers this space.
BAMSurgeon\autocite{ewing2015bamsurgeon} spikes variants into existing BAMs but uses deterministic VAF assignment.
ART\autocite{huang2012art} and NEAT\autocite{stephens2016neat} generate reads de novo but lack tumour models and variant injection.
ReSeq\autocite{schmeing2021reseq} achieves high read-level realism but supports no variant or cancer features.
No genome-wide UMI or duplex simulator exists.
No tool combines cfDNA fragmentation with somatic variant spike-in.

\varforge{} addresses these gaps in a single binary.
The contributions are:
\begin{enumerate}[leftmargin=*, label=(\roman*), nosep]
  \item A unified pipeline covering de-novo reads, variant injection, tumour modelling, UMI/duplex simulation, cfDNA fragmentation, and library artefacts, all via YAML configuration.
  \item The first genome-wide UMI and duplex sequencing simulator.
  \item The first in-silico cfDNA generator with somatic variant support.
  \item Stochastic VAF sampling via per-read Bernoulli trials.
  \item A streaming Rust implementation with bounded memory.
\end{enumerate}
